% ============================================================
% 204_Fraktale_Boundary condition_De_ch.tex
% Fraktale Boundary condition in den Torusmodefunktionen
% Perturbation theory in xi_0 erster order
% Johann Pascher — May 2026
% ============================================================

\section*{Abstract}

Die Mode functions $\phi_{n,l,j}$ aus Doc.~202 (Teil II) gelten
für einen einfachen Torus mit $D_f = 3$. Die fraktale Dimension
$D_f = 3-\xi_0$ des FFGFT-Torus führt zu einer correction erster
order in $\xi_0$. Diese correction modifiziert Normalisation und
Dispersion relation, lässt aber die Grundstruktur der Wellenfunktionen
unchanged. Die Massen entstehen weiterhin als Eigenvalue-Residuen
nach fast vollständiger Kompensation.

%------------------------------------------------------------
\section{Starting point: Einfacher Torus (\texorpdfstring{$D_f = 3$}{Df = 3})}
%------------------------------------------------------------

Auf dem einfachen Torus $T^3$ mit Seitenlänge $L_0$ gilt das
Standardintegrationsmaß:
\begin{equation}
  d\mu_0 = d^3x, \qquad
  \int_{T^3} d\mu_0 = L_0^3 = V
\end{equation}

Die Mode functions (Doc.~202):
\begin{equation}
  \phi_{n,l,j}^{(0)}(t,\mathbf{x})
  = \frac{1}{\sqrt{V}}\,
    e^{i\mathbf{k}_{n,l}\cdot\mathbf{x}} \cdot
    Y_{l,j}(\theta,\varphi) \cdot
    e^{-im_{n,l,j}t}
  \label{eq:phi0}
\end{equation}

sind exakt normiert: $\int_{T^3}|\phi^{(0)}|^2\,d\mu_0 = 1$.

%------------------------------------------------------------
\section{Das fraktale measure}
%------------------------------------------------------------

\subsection{Definition}

Auf dem FFGFT-Torus mit fraktaler Dimension $D_f = 3-\xi_0$ wird
das Integrationsmaß modifiziert. Für ein fraktales Medium mit
Hausdorff-Dimension $D_f$ gilt das skalierte measure:

\begin{equation}
  d\mu_{D_f} = \left(\frac{|\mathbf{x}|}{L_0}\right)^{D_f - 3} d^3x
  = \left(\frac{|\mathbf{x}|}{L_0}\right)^{-\xi_0} d^3x
  \label{eq:fraktales-mass}
\end{equation}

\subsection{Expansion für kleines \texorpdfstring{$\xi_0$}{xi0}}

Da $\xi_0 = 4/30000 \ll 1$, entwickeln wir:
\begin{equation}
  \left(\frac{|\mathbf{x}|}{L_0}\right)^{-\xi_0}
  = 1 - \xi_0 \ln\frac{|\mathbf{x}|}{L_0} + \mathcal{O}(\xi_0^2)
  \label{eq:entwicklung}
\end{equation}

Das fraktale measure ist also:
\begin{equation}
  d\mu_{D_f} = \left(1 - \xi_0\ln\frac{|\mathbf{x}|}{L_0}\right)d^3x
  + \mathcal{O}(\xi_0^2)
  \label{eq:mass-entwickelt}
\end{equation}

%------------------------------------------------------------
\section{Korrigierte Normalisation erster order}
%------------------------------------------------------------

\subsection{Das Normalisationsintegral}

Die correctede Normalisationsbedingung lautet:
\begin{equation}
  \int_{T^3}|\phi_{n,l,j}|^2\,d\mu_{D_f} = 1
\end{equation}

Mit dem Ansatz $\phi_{n,l,j} = N_{n,l,j} \cdot \phi^{(0)}_{n,l,j}$:
\begin{align}
  N_{n,l,j}^{-2}
  &= \int_{T^3}|\phi^{(0)}|^2\,d\mu_{D_f} \nonumber \\
  &= \int_{T^3}\frac{1}{V}\left(1 - \xi_0\ln\frac{|\mathbf{x}|}{L_0}\right)d^3x
  \label{eq:norm-integral}
\end{align}

\subsection{Calculation des correctionterms}

Das Integral des correctionterms:
\begin{equation}
  I_{\text{frak}} = -\frac{\xi_0}{V}
  \int_{T^3}\ln\frac{|\mathbf{x}|}{L_0}\,d^3x
  \label{eq:I-frak}
\end{equation}

Auf dem würfelförmigen Torus $[0,L_0]^3$ ist $|\mathbf{x}|$
der Abstand vom Ursprung. Das Integral liefert:
\begin{equation}
  \frac{1}{V}\int_{T^3}\ln\frac{|\mathbf{x}|}{L_0}\,d^3x
  = \ln\langle r\rangle/L_0
  \approx \ln(c_0)
\end{equation}

mit einer geometrischen Konstanten $c_0 \approx 0{,}5$ (mittlerer
relativer Abstand im Torus). Damit:
\begin{equation}
  N_{n,l,j}^{-2} = 1 + \xi_0|\ln c_0| + \mathcal{O}(\xi_0^2)
  \approx 1 + 0{,}69\,\xi_0
  \label{eq:N-quadrat}
\end{equation}

\subsection{Korrigierter Normalisationsfaktor}

Bis zur ersten order in $\xi_0$:
\begin{equation}
  N_{n,l,j} = 1 - \frac{0{,}69}{2}\,\xi_0 + \mathcal{O}(\xi_0^2)
  \approx 1 - 4{,}6 \times 10^{-5}
  \label{eq:N-wert}
\end{equation}

Die correction ist tiny ($\sim 10^{-5}$) — consistent damit,
dass $\xi_0 \ll 1$.

%------------------------------------------------------------
\section{Korrigierte Dispersion relation}
%------------------------------------------------------------

\subsection{Der fraktale Laplace operator}

Auf dem fraktalen Torus wird der Laplace operator modifiziert.
Für kleine $\xi_0$ gilt:
\begin{equation}
  \Delta_{D_f} = \Delta_3 + \xi_0 \cdot \delta\Delta + \mathcal{O}(\xi_0^2)
  \label{eq:frak-laplace}
\end{equation}

Der correctionoperator $\delta\Delta$ entsteht aus der
measuremodifikation:
\begin{equation}
  \delta\Delta\,\phi
  = -\nabla\!\left(\ln\frac{|\mathbf{x}|}{L_0}\right)\cdot\nabla\phi
  = -\frac{\mathbf{x}}{|\mathbf{x}|^2}\cdot\nabla\phi
  \label{eq:delta-Delta}
\end{equation}

\subsection{Eigenvaluekorrektur}

Die correctede Klein-Gordon-Gleichung:
\begin{equation}
  \left[\Box + \xi_0\cdot\delta\Box\right]\phi_{n,l,j}
  = -m_{n,l,j}^2\,\phi_{n,l,j}
\end{equation}

Für die ebene Welle $e^{i\mathbf{k}\cdot\mathbf{x}}$ liefert
$\delta\Delta$ in erster order:
\begin{equation}
  \langle\phi^{(0)}|\delta\Delta|\phi^{(0)}\rangle
  = -\langle\frac{\mathbf{x}}{|\mathbf{x}|^2}\cdot\nabla\rangle
  = -i\,k\langle\cos\theta_{kx}/r\rangle
\end{equation}

Auf dem Torus mittelt sich dieser Term zu null, da
$\langle \mathbf{x}/|\mathbf{x}|^2 \rangle = 0$ (Symmetrie).

Die eigenfrequency-correction kommt aus dem skalaren Teil:
\begin{equation}
  \delta\omega_{n,l,j}^2
  = -\xi_0 \cdot k_{n,l}^2 \cdot \langle\ln\frac{|\mathbf{x}|}{L_0}\rangle
  \approx 0{,}69\,\xi_0 \cdot k_{n,l}^2
  \label{eq:omega-korrektur}
\end{equation}

%------------------------------------------------------------
\section{Die correcteden Mode functions}
%------------------------------------------------------------

Die vollständigen Mode functions bis zur ersten order in $\xi_0$:
\begin{equation}
  \boxed{
    \phi_{n,l,j}^{(D_f)}(t,\mathbf{x})
    = \frac{N_{n,l,j}}{\sqrt{V}}\,
      e^{i\mathbf{k}_{n,l}\cdot\mathbf{x}} \cdot
      Y_{l,j}(\theta,\varphi) \cdot
      e^{-im_{n,l,j}^{(D_f)}\,t}
  }
  \label{eq:phi-corrected}
\end{equation}

mit:
\begin{align}
  N_{n,l,j} &= 1 - 4{,}6\times10^{-5} + \mathcal{O}(\xi_0^2)
  \label{eq:N-final} \\
  \left(m_{n,l,j}^{(D_f)}\right)^2
  &= k_{n,l}^2\left(1 + 0{,}69\,\xi_0\right) + \langle\hat{V}\rangle
  \label{eq:m-corrected}
\end{align}

Die fraktale correction $0{,}69\,\xi_0 \approx 3 \times 10^{-5}$
der kinetischen Energie ändert die Massenentstehung nicht qualitativ:
Die Massen bleiben Eigenvalue-Residuen nach fast vollständiger
Kompensation — die correction tritt nur in der 5. Nachkommastelle auf.

%------------------------------------------------------------
\section{Physikalische Significance}
%------------------------------------------------------------

\subsection{Warum die correction klein bleibt}

Die fraktale correction ist von order $\xi_0 \approx 1{,}3\times10^{-4}$.
Das ist consistent mit der mass hierarchy: Die Teilchenmassen sind
von order $\xi_0^{p_i} \cdot v$, also selbst durch $\xi_0$-Potenzen
suppressed. Die Wellfunktionskorrektur liegt in derselben order.

\subsection{Orthogonalität bleibt erhalten}

Da die correction \eqref{eq:N-final} mode-independent ist
(dasselbe $N$ für alle $(n,l,j)$), bleibt die Orthogonalität:
\begin{equation}
  \int_{T^3}
  \phi_{n,l,j}^{(D_f)*}\,\phi_{n',l',j'}^{(D_f)}\,d\mu_{D_f}
  = \delta_{nn'}\delta_{ll'}\delta_{jj'}
\end{equation}

Das Nichtabschluss-Theorem (Doc.~193) gilt unchanged.

\subsection{Completeness}

Das Massenentstehungs-Bild aus Doc.~202 (Teil V) bleibt korrekt:
\[
  m_i^2 = (2\pi n/L_0)^2\cdot(1+0{,}69\,\xi_0) + \langle\hat{V}\rangle
  \approx (2\pi n/L_0)^2 + \langle\hat{V}\rangle
\]

Die correction $(1+0{,}69\,\xi_0)$ ist numerisch irrelevant im Vergleich
zur Kompensation auf 32--39 Größenordnungen.

%------------------------------------------------------------
\section*{Summary}
%------------------------------------------------------------

Die fraktale Boundary condition aus $D_f = 3-\xi_0$ modifiziert die
Mode functions von Doc.~202 in erster order:
\begin{itemize}
  \item Normalisationsfaktor: $N = 1 - 4{,}6\times10^{-5}$
  \item Dispersionskorrektur: $+0{,}69\,\xi_0$ auf die kinetische Energie
\end{itemize}
Beide correctionen sind von order $\xi_0 \approx 1{,}3\times10^{-4}$
und ändern die physikalische Struktur nicht. Die Task aus Doc.~202,
Problem~2 ist damit formal abgeschlossen.


%------------------------------------------------------------
\section*{Notation}
%------------------------------------------------------------

\begin{center}
{\footnotesize
\resizebox{\textwidth}{!}{\begin{tabular}{lp{9cm}}
\toprule
\textbf{Symbol} & \textbf{Meaning} \\
\midrule
$D_f = 3-\xi_0$ & Fractal dimension of the FFGFT torus \\
$\xi_0 = 4/30000$ & Geometric base constant \\
$d\mu_0 = d^3x$ & Standard measure (simple torus, $D_f=3$) \\
$d\mu_{D_f}$ & Fractal measure ($D_f = 3-\xi_0$) \\
$\phi^{(0)}_{n,l,j}$ & Unperturbed mode functions ($D_f=3$) \\
$\phi^{(1)}_{n,l,j}$ & First-order perturbation correction in $\xi_0$ \\
$N_{n,l,j}$ & Normalisation factor (corrected) \\
$N \approx 1 - 4.6\times10^{-5}$ & Numerical value of normalisation correction \\
$\delta\Delta$ & Correction to the Laplacian from $D_f$ \\
$\delta\omega^2_{n,l,j}$ & First-order eigenfrequency correction \\
$0.69\,\xi_0$ & Dispersion correction ($\approx 3\times10^{-5}$) \\
$\phi^{(D_f)}_{n,l,j}$ & Corrected mode functions to $\mathcal{O}(\xi_0)$ \\
$m^{(D_f)}_{n,l,j}$ & Corrected mass to $\mathcal{O}(\xi_0)$ \\
\bottomrule
\end{tabular}}
}
\end{center}

\medskip\noindent
\textbf{References:} Doc.~180 ($\xi_0$), Doc.~193 (Nichtabschluss),
Doc.~202 (Mode functions, Massenentstehung).
